% !PRINTED 185
\weberpage{185}
%<<CONT
one obtains six relations among the nine coefficients
of $A$.

These relations read, if

\weq{3}{A = \begin{pmatrix} a_1, & a_2, & a_3\\ b_1, & b_2, & b_3\\ c_1, & c_2, & c_3 \end{pmatrix}}

is assumed:

\weq{4}{\begin{array}{r@{}c@{}l@{\qquad}r@{}c@{}l}
a_1^2+b_1^2+c_1^2 &{}={}& 1, & a_2a_3+b_2b_3+c_2c_3 &{}={}& 0\\
a_2^2+b_2^2+c_2^2 &{}={}& 1, & a_3a_1+b_3b_1+c_3c_1 &{}={}& 0\\
a_3^2+b_3^2+c_3^2 &{}={}& 1, & a_1a_2+b_1b_2+c_1c_2 &{}={}& 0,
\end{array}}

and are well known from analytic geometry. For the square
of the substitution determinant $|A|$ these relations yield, by
the multiplication rule of determinants, the value 1, and consequently
$|A|$ has the value $\pm 1$.

From the formulae (4) it follows that the \gesperrt{inverse substitution} of $A$

\[
A^{-1} = \begin{pmatrix} a_1, & b_1, & c_1\\ a_2, & b_2, & c_2\\ a_3, & b_3, & c_3 \end{pmatrix}
\]

is identical with the transposed one, and this property could
also serve as a definition of the orthogonal substitutions.

The totality of the orthogonal substitutions forms
a group. Among them there is contained a narrower group,
distinguished by the value

\weq{5}{|A| = + 1}

which we shall designate as the group of the \gesperrt{proper
orthogonal substitutions}.

One may interpret the linear substitution (1) as the transition
from one rectangular coordinate system to a second
with the same origin, if one regards $x_1, x_2, x_3$ and
$y_1, y_2, y_3$ as rectangular coordinates of one and the same
(variable) point in the first and in the second coordinate system. By $A$ the mutual position of the
two coordinate systems is determined, and conversely. If
the condition (5) holds, as we shall now assume, then
the first coordinate system can be brought into coincidence with the second by rotation about a fixed axis through a
definite angle of rotation.
